% Round 18. Root proof; separate independent review required.
\section{All codimensions in odd-dimensional product smoothings}
\label{r18od:section}

The middle numerical groups of even-dimensional complete intersections
need not be generated by the hyperplane class.  This section retains
those groups.  In the present model their extra part is killed by
the literal Gysin map, and chosen component representatives have
restrictions in an explicitly nondegenerate subspace.  These two
facts replace any unproved middle cycle comparison or Hodge-standard
assumption.

\subsection{Geometric input and nonmiddle complete-intersection classes}

Use exactly the hypotheses and the ordered resolved model of NF18,
including all smoothness, transversality, connectedness and generic
smoothness assumptions.  Suppose now that
\[
                         n=2k+1\ge5.
 \tag{OD18.1}\label{r18od:dimension}
\]
Linear factors are included.  Write $h$ for the pullback of the
ambient hyperplane class on the whole model.  The geometric
description established in HP18/NF18 gives:
\begin{itemize}
\item each component $D_i$ is a smooth complete intersection
of dimension $2k+1$, followed by the specified codimension-two
blowups in smooth complete intersections of dimension $2k-1$;
\item each double $E_{ij}$ starts from a smooth complete intersection
of dimension $2k$, followed by the specified codimension-two
blowups in smooth complete intersections of dimension $2k-2$;
\item every center in these component and double-stratum blowup
sequences is geometrically integral, and its strict transform
under earlier steps is isomorphic to that original center, with
the ambient class $h$ identified with its hyperplane class.
\end{itemize}
The two blowup sequences are the actual native sequences of NF18;
no factorization of an arbitrary birational map is invoked.

We recall the precise comparison used for these varieties.
If $Y$ is a smooth geometrically integral complete intersection
of dimension $b$ and $2u\ne b$, its geometric cohomology in
degree $2u$ is $\mathbf Q_\ell h^u$.  For an actual native
rational cycle $z$ of codimension $u$, the coefficient of its
geometric class in that line is
\[
           \frac{\deg_Y(z h^{b-u})}{\deg_Y(h^b)}\in\mathbf Q.
 \tag{OD18.2}\label{r18od:rational-coefficient}
\]
Subtracting that rational multiple of the actual $h^u$ gives a
homologically zero, hence numerically zero, cycle.  Conversely a
numerically zero cycle has coefficient zero in this formula.
It follows, for precisely these nonmiddle degrees, that
$N^u(Y)_{\mathbf Q}=\mathbf Qh^u$ and that numerical triviality
of actual cycles implies geometric homological triviality.
This is a proved comparison in the stated degree, not an assumed
cycle-class map on an arbitrary numerical quotient.
For odd $b$, all algebraic cycle degrees are nonmiddle and the
argument applies to every $u$.  The geometric cohomology input
is the smooth complete-intersection weak Lefschetz theorem
(Ghorpade--Lachaud, arXiv:0808.2169, Propositions 3.1--3.2),
together with Poincare duality and the nonzero hyperplane degree.

\subsection{Component bases and the one exceptional degree on the doubles}

For a component let $\Gamma_{i,j}$ denote the total-transform
exceptional divisor associated with its center indexed by $j>i$.
The numerical blowup decomposition and
\eqref{r18od:rational-coefficient} give, for every
$1\le u\le n-1$,
\[
  N^u(D_i)_{\mathbf Q}
     =\mathbf Qh_i^u
       \oplus\bigoplus_{j>i}\mathbf Q\,\Gamma_{i,j}h_i^{u-1}.
 \tag{OD18.3}\label{r18od:component-basis}
\]
Empty centers are omitted.  All entries in this basis have actual
native Chow representatives as displayed.  Indeed a codimension-two
blowup adds $j_*p^*N^{u-1}(Z)$, and $h|_Z$ is the specified
hyperplane class; its generator is exactly the exceptional divisor
times $h^{u-1}$.  Later pullbacks give total transforms.  The same
Chow decomposition and geometric cohomology blowup formula show
that these cycle classes span the entire geometric even cohomology
of the component in the indicated degree.  No assertion about
its odd cohomology is made or needed.

For an edge $e=ij$, put $t_e=\#\{l>j:Z_{ij,l}\ne\varnothing\}$
and let $f_{e,l}$ be its total-transform exceptional divisors.
For $1\le u\le 2k-1$ define the native numerical subspace
\[
 T_e^u=\operatorname{span}_{\mathbf Q}
       \{h_e^u,\ f_{e,l}h_e^{u-1}:l>j\}
             \ \subseteq\ N^u(E_e)_{\mathbf Q}.
 \tag{OD18.4}\label{r18od:edge-tautological}
\]
These displayed generators are independent: this follows from
the numerical blowup decomposition and the nonzero degree of
every hyperplane power on each center, including its middle
power.  In particular $\dim T_e^u=1+t_e$ in every indicated degree.
For $u\ne k$, equation \eqref{r18od:rational-coefficient}, applied
to the base and to each center in the blowup decomposition, gives
\[
                          N^u(E_e)_{\mathbf Q}=T_e^u.
 \tag{OD18.5}\label{r18od:only-middle-extra}
\]
There may be additional classes in $N^k(E_e)$.

\begin{lemma}[The middle subspace and its orthogonal complement]
\label{r18od:middle-splitting}
The middle intersection pairing on $T_e^k$ is nondegenerate.
In the total-transform basis of \eqref{r18od:edge-tautological}
it has the diagonal entries
\[
       \deg_{H_i\cap H_j}(h^{2k}),\qquad
                      -\deg_{Z_{ij,l}}(h^{2k-2})\quad(l>j).
 \tag{OD18.6}\label{r18od:middle-matrix}
\]
Consequently there is a canonical rational orthogonal splitting
\[
     N^k(E_e)_{\mathbf Q}=T_e^k\mathbin{\perp}P_e^k.
 \tag{OD18.7}\label{r18od:middle-orthogonal}
\]
Here $P_e^k$ is allowed to be nonzero and has no asserted
positivity or numerical-to-homological comparison.
\end{lemma}

\begin{proof}
For one codimension-two blowup $q:\widetilde Y\to Y$, put
$j:E_Z\hookrightarrow\widetilde Y$, $p:E_Z\to Z$ and
$\xi=c_1(\mathcal O_{E_Z}(1))$.  At middle degree $k$ the
blowup decomposition is $q^*N^k(Y)\oplus j_*p^*N^{k-1}(Z)$.
Projection gives zero cross pairing, since $p_*1=0$.
Self-intersection $j^*j_*=-\xi$ and $p_*\xi=1$ give the
exceptional pairing $-\deg_Z(\alpha\beta)$.
Apply this successively, pulling all earlier summands back.
The original hyperplane line has square $\deg(h^{2k})>0$;
each displayed exceptional hyperplane line has the negative
nonzero square in \eqref{r18od:middle-matrix}.
This proves the matrix and its nondegeneracy.  The full middle
numerical pairing is nondegenerate by numerical duality, so
its rational orthogonal complement gives the asserted splitting.
Every degree is the native degree on the displayed proper
variety; no degree on the absolute total-model Chow group occurs.
\end{proof}

Put $T^u=\bigoplus_eT_e^u$ and $P^k=\bigoplus_eP_e^k$.
In other interior degrees put $P^u=0$, so that
\[
                              W^u=T^u\oplus P^u.
 \tag{OD18.8}\label{r18od:W-splitting}
\]
Let $\operatorname{pr}_T^u$ be the rational projection from this
splitting; it is the identity away from the middle degree.

\subsection{Actual restrictions and the literal kernel}

\begin{lemma}[Restriction images are in the displayed subspaces]
\label{r18od:restrictions}
For every $1\le u\le n-1$ for which the edge degree is interior,
\[
                                \rho_u(C^u)\subseteq T^u.
 \tag{OD18.9}\label{r18od:rho-in-T}
\]
Moreover, choose the actual basis representatives in
\eqref{r18od:component-basis} for a component numerical tuple.
Their geometric restrictions, and their differences on a double,
belong to the geometric cohomology span of the corresponding
actual generators in \eqref{r18od:edge-tautological}.
In particular,
\[
                                  P^k\subseteq\ker\sigma_{k+1}.
 \tag{OD18.10}\label{r18od:primitive-killed}
\]
\end{lemma}

\begin{proof}
The restriction of $h_i^u$ is $h_e^u$.  The restriction of
$\Gamma_{i,j}h_i^{u-1}$ is the actual divisor restriction
$\Gamma_{i,j}|_{E_e}$ multiplied by the same global $h_e^{u-1}$.
The whole numerical divisor group of $E_e$ has basis
$h_e,f_{e,l}$, since $k\ge2$ and degree one is nonmiddle on
its base, while a codimension-two blowup adds only $N^0$ of
each connected center in divisor degree.  Thus the numerical
restriction is a rational linear combination of the displayed
degree-$u$ generators, proving \eqref{r18od:rho-in-T}.

This divisor expansion is valid also for geometric cohomology:
the base has $H^2=\mathbf Q_\ell h$, and the blowup formula
adds one geometric $H^0$ class for each geometrically integral
center.  The displayed native divisors span these classes.
The coefficients are the same rational numerical coefficients,
as is seen by pairing with a rational dual basis in numerical
degree $2k-1$; that degree has the nonmiddle comparison already
proved in \eqref{r18od:only-middle-extra}.  Hence the difference
between the actual divisor and its displayed rational expansion
has geometric class zero.  Multiplication by the actual global
$h^{u-1}$ gives the stated claim for the chosen representatives.
This does not assert a Chow equality for an arbitrary divisor
or a cycle-class map on $N^k(E_e)$.

Finally take $p\in P^k$ and $v\in C^k$.  The native signed
adjunction of DQ17 gives
\[
                  \langle\sigma_{k+1}p,v\rangle_C
                         =-\langle p,\rho_kv\rangle_W=0,
\]
by \eqref{r18od:rho-in-T} and orthogonality in
\eqref{r18od:middle-orthogonal}.  The pairing
$C^{k+1}\times C^k$ is perfect because $n=2k+1$.
Therefore $\sigma_{k+1}p=0$, proving
\eqref{r18od:primitive-killed}.  No cycle class for $p$ was used.
\end{proof}

\subsection{One global polarization and the weaker induction hypothesis}

Choose one relatively ample line bundle $\mathcal L$ on the
resolved model and put
$H_M=c_1(\mathcal L)+Mh$ for an integer $M\ge0$.
It is relatively ample for all sufficiently large $M$.
The ambient class $h$ itself sends each basis vector of $T_e^u$
to its identically indexed basis vector of $T_e^{u+1}$ for
$1\le u\le 2k-2=n-3$.  Consequently
\[
 \operatorname{pr}_T^{u+1}\circ H_M|_{T^u}
                              =MI+C_u
 \qquad(1\le u\le n-3)
 \tag{OD18.11}\label{r18od:finite-matrices}
\]
in these bases, for fixed rational square matrices $C_u$.
There are finitely many of them, and each determinant is monic.
Choose one sufficiently large integer $M$ avoiding all their
roots, and write $H=H_M$.  It follows that
\[
 \begin{gathered}
 H|_{T^u}\text{ is injective},\\
 W^{u+1}=H T^u+P^{u+1}
                       \qquad(1\le u\le n-3).
 \end{gathered}
 \tag{OD18.12}\label{r18od:ample-consequences}
\]
No assertion that $H$ preserves the subspaces $T^u$ is needed:
the projection in \eqref{r18od:finite-matrices} is intentional.
The extra summand in the second line is zero except at degree
$k$, and there it lies in the literal Gysin kernel by
\eqref{r18od:primitive-killed}.

\begin{lemma}[Induction with injectivity only on the restriction image]
\label{r18od:weaker-DG}
On any reduced strictly semistable model, let $H$ be one global
divisor class and $r\ge3$.  Suppose that
\[
 \begin{gathered}
 U_{r-1}=0,\qquad
 W^{r-1}=H W^{r-2}+\ker\sigma_r,\\
 H\text{ is injective on }
              \operatorname{im}(\rho_{r-1}\sigma_{r-1})
                      \subseteq W^{r-1}.
 \end{gathered}
 \tag{OD18.13}\label{r18od:weak-hypotheses}
\]
Then $U_r=0$.
\end{lemma}

\begin{proof}
Represent a numerator class by $z=\sigma_rw$ with $\rho_rz=0$.
Write $w=Hb+t$ with $t\in\ker\sigma_r$.  Since restriction,
Gysin and $A$ commute with multiplication by one global class,
\[
          z=H\sigma_{r-1}b,
        \qquad 0=\rho_rz=H\rho_{r-1}\sigma_{r-1}b.
\]
The stated injectivity implies
$\rho_{r-1}\sigma_{r-1}b=0$.
By $U_{r-1}=0$ there is $v\in C^{r-2}$ with
$\sigma_{r-1}b=A_{r-2}v$.  Therefore
$z=HA_{r-2}v=A_{r-1}(Hv)$, as required.
This proves the assertion on the displayed rational numerical
groups without an injectivity assertion on all of $W^{r-1}$.
\end{proof}

\begin{proposition}[All middle numerical defects vanish]
\label{r18od:all-U}
For the model of \eqref{r18od:dimension},
\[
                               U_r=0
                     \qquad(2\le r\le n-2).
 \tag{OD18.14}\label{r18od:U-vanishes}
\]
\end{proposition}

\begin{proof}
The initial assertion $U_2=0$ is HG18 for linear factors and
NF18 for general factors.  For each $3\le r\le n-2$, the
second line of \eqref{r18od:ample-consequences}, with $u=r-2$,
and \eqref{r18od:primitive-killed} give
$W^{r-1}=H W^{r-2}+\ker\sigma_r$.
The image $\rho_{r-1}\sigma_{r-1}(W^{r-2})$ lies in
$T^{r-1}$ by \eqref{r18od:rho-in-T}.
The first line of \eqref{r18od:ample-consequences}, now with
$u=r-1\le n-3$, gives the injectivity required in
Lemma~\ref{r18od:weaker-DG}.  Apply that lemma successively.
The same single ample class works in every degree.
\end{proof}

\subsection{Chosen representatives close the remaining compatibility gate}

\begin{proposition}[Compatibility without a general middle cycle comparison]
\label{r18od:K-equals-B}
For every $2\le a\le n-2$ one has
\[
                     K_{\mathrm{comp},\ell}^{\,a}=\ker\rho_a.
 \tag{OD18.15}\label{r18od:K-B}
\]
\end{proposition}

\begin{proof}
The inclusion from left to right is frozen.  Let a native
numerical tuple $v$ lie in $\ker\rho_a$.  On each component
choose the rational actual representative obtained by its
unique expansion in \eqref{r18od:component-basis}.
By Lemma~\ref{r18od:restrictions}, the geometric difference
of these chosen restrictions on each double is a rational
linear combination of the classes of its displayed $T_e^a$
generators.

If $a\ne k$, the same numerical blowup decomposition and
nonmiddle complete-intersection comparison used in
\eqref{r18od:only-middle-extra} show that a numerically zero
actual edge cycle in that degree has geometric class zero.
One may also check this by pairing the geometric expansion
with the complementary native hyperplane and exceptional
classes, whose rational intersection matrix is perfect.

If $a=k$, use the explicit diagonal matrix
\eqref{r18od:middle-matrix}.  The actual restriction difference
is numerically zero because $\rho_a v=0$.  Pair its geometric
expansion with each actual $T_e^k$ generator.  All intersection
numbers are zero, and the displayed diagonal matrix is
invertible over $\mathbf Q$.  Every coefficient in the
geometric expansion is therefore zero.  This proves geometric
compatibility on that double too.

The chosen native rational representatives are thus geometrically
compatible on every double and exhibit $v\in K_{\mathrm{comp},\ell}^a$.
Only this chosen representative family was used.  No assertion
that an arbitrary numerically zero middle cycle on an even
dimensional complete intersection is homologically zero has
been introduced, and no Chow equality of arbitrary compatible
restrictions has been claimed.
\end{proof}

\begin{theorem}[Full local strong G in all odd dimensions covered by the model]
\label{r18od:full-G}
Under the exact NF18 product-smoothing hypotheses with
$n=2k+1\ge5$, every specified geometrically homologically
trivial rational cycle on the generic fiber satisfies strong G
in every codimension.  In particular this covers the HP18
hyperplane subfamily.  The conclusion reapplies to all cycles
over every unramified finite valued extension.
\end{theorem}

\begin{proof}
For $2\le r\le n-2$, put $a=n-r$.  Proposition
\ref{r18od:all-U} gives $U_r=0$, and Proposition
\ref{r18od:K-equals-B} gives $K_a=B_a$.
The exact sequence DQ17 therefore makes
$V_a/(K_a+L_a)$ zero in every such degree.
ER17 and VT17 give obstruction zero; the frozen reconstruction
then supplies the required ordinary actual form on the exact
complement of the original specified cycle, against every
compact ordinary actual test.  The frozen divisor, one-cycle
and zero-cycle theorems cover the remaining degrees.
An unramified extension preserves the model hypotheses;
apply the proof afresh to each cycle over the new field.
No claim that new cycles descend as pullbacks is involved.
\end{proof}

\paragraph{Consequences and limits.}
Together with EV18, this supplies full local G in every
dimension $n\ge4$ for the entire stated linear and nonlinear
product-smoothing model class.  The global GF18/NG18
Chinese-remainder construction then has full coverage in all
codimensions, whenever this local theorem is used at the
distinguished prime and the frozen GP17 coverage is used at
every other finite place.  Its complete complementary
Beilinson--Bloch height pairings follow from PH16 and GA17.
The restriction to unramified extensions at that distinguished
prime remains in force unless a separate ramification theorem
is proved.

The proof does not turn a general strictly semistable model
into this one.  For a general model the component restrictions
need not lie in a subspace such as $T$, and the orthogonal
middle complement need not lie in the Gysin kernel.
General $U_r$ and the exact homological period-vanishing
problem therefore remain open.  All conclusions are relative
to the final frozen analytic and height interfaces; no
historical-priority or external-refereeing claim is made.
